\subsection{Academic Achievements}
Figure [] shows the extent of grade inflation among students according to their previous academic achievements. I find that students who scored higher in Mathematics in GCSE receive top grades with higher probabilities. A student with a GCSE score of 8 in mathematics receives a top grade by 11 percentage points more than in the standardized test regime. The probability increments are smaller for students with lower grades, with a 1 unit decrease in the score accompanied by an approximately 4 percentage points reduction in their probability of obtaining A or A*. Students with a math score of 5 receives top grades with equal probability as the standardized exams, and students with grade 4 in reduce their chances of receiving a top grade by 5 percentage points. The patterns are similar between 2020 and 2021 but with the levels in the increments being slightly higher for all grades in 2021. 

Figures [] plot the marginal treatment effects after controlling for subject composition. The increments across grade points are similar to figure [], but with difference in levels. In 2020, a student with a score of math at 4 or 5 reduces their chances achieving A or A*, by 12 and 6 percentage points respectively. In 2021, students with a math score of 4 achieves A or A* with less probability than in the standardized regime.

Figure [] plots the resulting probability of scoring a top grade across math grades for both standardized and non-standardized grading regimes. In both 2020 and 2021 effect of a higher grade in GCSE, the standardized math exam taken at the beginning of secondary schools, on the probability of achieving a top grade is larger in the non-standardized regime than in the standardized regime. The average increments in the probability of achieving A or A* per a increase in the GCSE grade in approximately 6 percentage points in the standardized grading regime, but the increments for the non-standardized regime is 7 percentage points. The contrast between the two edges of the GCSE score distribution stands out, as the probability difference in achieving A or A* between a GCSE score of 8 and 4 is 50 percentage points for the non-standardized testing regime, and the gap difference is 37 percentage points for the standardized testing regime.

Figures [] plot the marginal treatment effects after removing the school level controls. The removal of the school level controls suppresses the achievements predicted by the students GCSE scores. The probability difference between a score of 8 and 4 in math in GCSE is 11 percentage point after removing the school level controls, and 0 percentage points after removing the school level and demographic control all together. The the gap is statistically smaller than the gap with the specification with the school level controls, which was 17 percentage points.

The results indicate that academic achievements are still reflected into the students chances of achieving a top grade in the non-standardized regime. Teachers or school still refer to their students academic performance in the past and grade students according to their records. As the increments by grades are constant between figure [] and figure[], student rank across previous academic achievement levels are preserved within classrooms, as most subject courses are provided in a single classroom. Moreover, the non-standardized grading regime expanded the achievement gaps across students with contrasting previous academic achievements. This implies that schools and teachers amplifies the achievement gaps between students with varying previous academic achievements.

I hypothesize that although teachers were not obligated to refer to students' academic achievements to decide their final grades, the teachers preserved the ordering of the students according to academic achievements. Teachers may have punished students with lower academic scores more severely than higher performing students, as lower performing student's grades were shifted upward due to the loosening of grading standards and teachers attempted to maintain the integrity of the grade scale by lowering the grades of student with the largest shifts in their final grades. 

However, the widened difference in final academic achievements across previous academic performance is dampened, as schools attended by pupils with lower academic achievements assign top grades with higher frequency to their pupil than schools attended by students with higher academic achievements.  As discussed in section [], schools with low value-added assigned top grades to students with higher tendency than schools with high value added. Accouting for students sorting patterns across schools with varying value added based on their academic achievements weakens the pronounced relation between previous academic performance and final academic achievements. Furthermore, the increments in increments across academic performance disappear after accounting for the demographic sorting of students to schools.  

\subsection{Demographics}
Table [] denotes results on grade assignment patterns across gender groups. I find female students were 6 percentage points more likely to receive a top grade than male students. The difference across gender groups exists in similar magnitudes for both 2020 and 2021. Column 2 and 3, which shows the result for 2020 and 2021 cohorts respectively, denotes a 6.3 and 5.2 percentage difference between male and female groups. The difference across gender groups remains after removing the school level controls. Column 4 reports a 5.7 percentage difference between female and male students. The magnitudes of the differences is almost equivalent to the impact of scoring an additional score in GCSE exams, which predicted a 6 percentage point increase in achieving a top grade for a unit increment in scores. 

The results indicate student's gender is reflected into their final grades, and shrinking the attainment gap between female and male students. The results indicates a re-ordering of students within a classroom group based on the students gender. The achievement gap between male and female students shrinks, leading to reduce the gender gap in achievements. I provide multiple 
hypothesize to explain this pattern; (1) difference in academic performance in the classroom projecting itself to the final grades, and (2) the population of female teachers among the teacher population. Female students tend to perform worse than male students in high stake exams, and research attributes its causes to anxiety (Ballen et al 2025). As the non-standardized A-levels did not include any in-person exams, the higher outcomes of female students in academic tasks in their classrooms may have been reflected in the final grades. Another explanation is related to the gender of teachers. As majority of secondary school teachers are female, teachers may have assigned better grades to students that shares the same gender as themselves.

Figure []-[] plots the grading patterns across parental occupation classes. The increments in probability points are similar across all occupation class. The largest increments in the probability of receiving a top grade is the students with parents with an occupation as [] with [] percentage points. The lowest was for students with parental occupation as [] with [] percentage points. However, students with parents in lower supervisory occupations (SES code 5) was the largest increments in both 2020 and 2021. Figure [] plots the resulting probability in achieving a top grade across parental occupation class for both non-standardized and standardized exams. Results indicates that the patterns across parents occupation classes and grades remain similar between the two grading regimes. The largest difference across occupational groups is the difference between 2 and 6 for both 2020 and with 9 percentage points and 6 percentage points respectively. The largest gap also between 2 and 6 for the standardized regime as well, with nearly 9 percentage point difference between the two classes. 

Figure [] - [] reports the difference in students probability on receiving a top grade using different measures of economic affluence. For all measures, the probability of receiving an A/A* is positively related with the students background wealth. Figure [] reports the results using the level of economic deprivation for the students residential address. A unit increase in the quintile of the deprivation score, which is descending with degree of wealth, is associated with a 1.6 percentage point increase in the students probability of achieving a top grade. I observe a 8 percentage point difference between students in the 1st and 5th quintile band in the income deprivation scores. Figure [] reports the difference in the achievement of high grades between regions with varying higher education progression rates. An one unit increase in the progression score (polar 4 score), which is ascending in the share of population progressing to higher education,  is associated with a 0.6 percentage point increase in the students probability of achieving a top grade. The gap between the lowest and higher progression score district is 4 percentage points. 

The results indicate that the demographics of the students and the affluence of their backgrounds influence the grades they receive. Although the patterns hint that the nonstandardized scheme worsened the achievement inequality across students with different backgrounds, the levels are subtle compared to student academic achievements. I hypothesize that affluence in students' backgrounds may influence the school/teachers perception of the students academic performance. Although teachers attempt not to use student's backgrounds into their grading, a bias towards rewarding more affluent students persists and slightly exacerbating the achievement inequality. 

[Needs fixing]I find that white students increased their odds of receiving high grades than non-white students. The odds for white students increased by 18\% than before, as non-white students improved their odds by 25\%. The column [] in Table [] denotes how the estimates barely change after removing the school-fixed effect interaction term from the model. The improvement in odds for white students was 18\%, which is identical to the estimate with the interaction term between school-fixed effect and treatment dummy. 



\subsection{Subject Type}
Table [] shows how grade inflation was concentrated in the subject groups. Column 1 denotes how facilitating subject, a class of subject that applicants are recommended take for applying to elite universities, were 8 percentage points less likely to be assigned a top grade than non-facilitating subjects. Columns 2 and 3 report the decomposition of the difference between 2020 and 2021, and the estimates are similar in magnitude. Columns 4 and 5 report the estimates after controlling for the quantitative content of the subject. The estimated difference between facilitating and nonfacilitating subjects decreases to 6 percentage points for both 2020 and 2021. Column 5 reports the estimates after removing the interaction terms between the school fixed effect and the treatment year dummy. The magnitudes of the estimate have slightly reduced to 0.77.

The results underscore how the content of the subjects influences the grading patterns. Subjects with more nonquantitative content were assigned a top grade with higher probability than subjects with less quantitative content. I attribute the difference to the difficulty in assigning a single grade to subjects with less quantitative content. Grades for subjects such as drama studies or media studies entrust a large degree of discretion on the graders, as the solutions of the contents are intrinsically difficult to measure. Teachers may feel morally guilty of ranking students in a strict order when they cannot confidently measure their outcomes, consequently narrowing the within-cohort grade distribution. In contrast, subjects with more quantitative content, such as mathematics or physics, are often easier to measure to a student's performance as solutions tend to be clear.

Table [] reports the regression coefficient between the increments in grades between 2020 and 2021 between and within the subject groups. The regression results are shown using the facilitation subjects as the subject group class. Columns 1 and 2 report the regression results using grade increments for facilitating subjects and nonfacilitating subjects, respectively, using an OLS regression. I find a serial correlation between subject groups within a school between the two years. For facilitating subjects, the average elasticity of grade increments is 3 percentage points for grade increments in facilitating subjects at 2020, and 1.2 percentage points for grade increments in nonfacilitating subjects at 2020. For nonfacilitating subjects, the average elasticity of grade increments is 3.5 percentage points for grade increments in non-facilitating subjects at 2020, and 0.9 percentage points for grade increments in facilitating subjects at 2020. Columns 3 and 4 report the coefficients using a median regression, to address the impact of outliers in the estimates of $\delta$s. The relative magnitudes of the coefficients are similar to the regressions using OLS. 

Figures [] and [] report the patterns among the subject groups in 2020 and 2021. I find a positive relationship between the same subject groups within a school. The 1 percentage point increase in facilitating subjects in 2020 is associated with an average [] percentage point increase facilitating subjects in 2021. 

My results demonstrate the dynamics of the grading decisions, as well as the degree of teacher discretion in their grading decisions. Firstly, the patterns highlight how schools react quickly to their expectations on how universities can strengthen their grading standards. As many universities were overloaded with inflow of students in 2020 due to the sharp rise in their applicants grades and grade requirements in their acceptance offers set before the pandemic erupted. I hypothesize that schools adjust the grading levels based on the expectations of the grade requirements so that their pupils do not be denied entry into their preferred institutions. Secondly, I infer that much of the grading decisions depended on the teachers and not the schools. Figures show that grade increments between 2020 and 2021 across subject groups do not show strong correlation, indicating that the grading decisions were centered around the subject levels, and not the school level. Lastly, the results point at the hesitance of teachers to tighten their grading standards after loosening them in the previous year. Most schools did not lower their grade increments in 2021 relative to their levels in 2020, while schools with lower degrees of grade increments in 2020 vastly increased their grade increment levels in 2021. The patterns highlight how teachers may find it difficult to tighten grading standards in contrast to loosening their standards. 